%-------------------------
% Resume in Latex
% Author: Reza Aliasgari Renani (reordered to reverse-chronological)
% License: MIT
%------------------------

\documentclass[a4paper, 10pt]{article} % Changed to 10pt to accommodate longer Russian text
\usepackage{titlesec}
\usepackage{marvosym}
\usepackage[usenames,dvipsnames]{color}
\usepackage{enumitem}
\usepackage[hidelinks]{hyperref}
\usepackage{fancyhdr}
\usepackage[T2A]{fontenc}
\usepackage[utf8]{inputenc}
\usepackage[noconfigs,russian,english]{babel}
\usepackage{tabularx}
\usepackage{fontawesome5}
\usepackage{multicol}
\usepackage[version=4]{mhchem}
\usepackage[
  a4paper,
  top    = 1.50cm, % Slightly reduced margins to fit content
  bottom = 1.50cm,
  left   = 1.70cm,
  right  = 1.70cm
]{geometry}
\usepackage{cmbright} % Sans-serif font family (Computer Modern Bright) for easier readability, supports Cyrillic with T2A

\renewcommand{\baselinestretch}{1.05} % Reduced stretch to fit more content
\renewcommand{\headrulewidth}{0pt}
\renewcommand{\footrulewidth}{0pt}

% Defining custom commands
\newcommand{\resumeItem}[1]{
    \item\small{
        {#1 \vspace{-2pt}}
    }
}

\newcommand{\resumeSubheading}[4]{
    \vspace{-2pt}\item
    \begin{tabular*}{1.0\textwidth}[t]{l@{\extracolsep{\fill}}r}
        \textbf{#1} & \textbf{\small #2} \\
        \small#3 & \small #4
    \end{tabular*}\vspace{-7pt}
}

\newcommand{\resumeProjectHeading}[2]{
    \item
    \begin{tabular*}{1.001\textwidth}{l@{\extracolsep{\fill}}r}
        \small#1 & \textbf{\small #2}\\
    \end{tabular*}\vspace{-7pt}
}

\renewcommand\labelitemi{$\vcenter{\hbox{\tiny$\bullet$}}$}
\renewcommand\labelitemii{$\vcenter{\hbox{\tiny$\bullet$}}$}

\newcommand{\resumeSubHeadingListStart}{\begin{itemize}[leftmargin=0.0in, label={}]}
\newcommand{\resumeSubHeadingListEnd}{\end{itemize}}
\newcommand{\resumeItemListStart}{\begin{itemize}}
\newcommand{\resumeItemListEnd}{\end{itemize}\vspace{-5pt}}

\setlength{\multicolsep}{-3.0pt}
\setlength{\columnsep}{-1pt}

\hypersetup{
    colorlinks=true,
    linkcolor=RoyalBlue,
    urlcolor=RoyalBlue,
    citecolor=RoyalBlue,
    anchorcolor=RoyalBlue
}

\input{glyphtounicode}

\pagestyle{fancy}
\fancyhf{} % clear all header and footer fields
\fancyfoot{}
\setlength{\footskip}{30pt} % Fixes \footskip warning
\urlstyle{same}
\raggedbottom
\raggedright
\setlength{\tabcolsep}{0in}

% Formatting section titles
\titleformat{\section}{
    \vspace{-4pt}\scshape\raggedright\large\bfseries
}{}{0em}{}[\color{black}\titlerule \vspace{-5pt}]

\pdfgentounicode=1

% Starting the document
\begin{document}

%----------HEADING----------
\begin{center}
    {\Huge \scshape {\fontsize{25}{30}\selectfont{Реза}} {\fontsize{25}{30}\selectfont{\textbf{Алиасгари Ренани}}}} \\ \vspace{3pt}
    \small ~ 
    \href{mailto:rezaaliasgarirenani@gmail.com}{\raisebox{-0.2\height}\faEnvelope\ \underline{rezaaliasgarirenani@gmail.com}} $|$ ~
    \href{https://orcid.org/0009-0000-8983-755X}{\raisebox{-0.2\height}\faOrcid\ \underline{ORCID}} $|$ ~
    \href{https://www.researchgate.net/profile/Reza-Aliasgari-Renani}{\raisebox{-0.2\height}\faResearchgate\
    \underline{ResearchGate}} $|$ ~
    \href{https://scholar.google.com/citations?user=L9Vv3C8AAAAJ&hl=en}{\raisebox{-0.2\height}\faGraduationCap\ \underline{Google Scholar}} $|$ ~
    \href{https://github.com/rezaaliasgarirenani}{\raisebox{-0.2\height}\faGithub\ \underline{GitHub}} $|$ ~
    \href{https://linkedin.com/in/reza-aliasgari-renani}{\raisebox{-0.2\height}\faLinkedin\ \underline{LinkedIn}}
    \vspace{-10pt}
\end{center}

%-----------EDUCATION-----------
\section{\textcolor{red}{Об}разование}
\resumeSubHeadingListStart
    %--- Most recent (M.Sc.) first ---
    \resumeSubheading
        {\href{https://mipt.ru/}{\parbox[t]{\dimexpr\linewidth-6cm}{Московский физико-технический институт (МФТИ, Физтех)}}}{Сентябрь 2024 -- Июнь 2026}
        {\textbf{Магистр} прикладной математики и физики, Программа: Физика плазмы, Средний балл: 4.7/5.0}{Москва, Россия}
        {\resumeItem{\textbf{Тема диссертации:} Исследование эффектов, вызванных радиацией, на системах обработки сигналов на основе FPGA для космических применений.}}
    %--- Earlier (B.Sc.) after ---
    \resumeSubheading
        {\href{https://mipt.ru/}{\parbox[t]{\dimexpr\linewidth-6cm}{Московский физико-технический институт (МФТИ, Физтех)}}}{Сентябрь 2020 -- Июнь 2024}
        {\textbf{Бакалавр} технической физики, Программа: Аэрокосмические технологии, Средний балл: 4.56/5.0}{Москва, Россия}
        {\resumeItem{\textbf{Тема диссертации:} Исследование эффектов облучения низкоэнергетичными (1 - 20 кэВ) электронами и высокоэнергетичными (1 МэВ) гамма-квантами на электрофизические свойства диэлектрик-полупроводниковых структур.}}
\resumeSubHeadingListEnd
\vspace{-10pt} % Adjusted spacing

%-----------EXPERIENCE-----------
\section{\textcolor{red}{Оп}ыт научных исследований}
\resumeSubHeadingListStart
    %--- Current / Most recent position first ---
    \resumeSubheading
        {\href{http://ai.mipt.ru/design-center}{\parbox[t]{\dimexpr\linewidth-6cm}{Дизайн-центр по проектированию микропроцессорной техники для систем с ИИ, Лаборатория разработки систем на кристалле}}}
        {Сентябрь 2024 -- Настоящее время}
        {\textbf{Программист / Инженер проектирования RTL}}{Москва, Россия}
        \resumeItemListStart
            \resumeItem{\textbf{Реализация DSP} \\ Перенос математических алгоритмов в эффективные реализации на Verilog. Создание библиотеки фиксированной точки и аппроксимаций функций на основе LUT (метод Горнера) для поддержки вычислений с фиксированной точкой. Реализация алгоритмов обработки изображений (rgb2hsv, сегментация цвета, обнаружение краев Собеля, глобальное тональное картирование, демозаицинг, вид с высоты птичьего полета, коррекция «рыбьего глаза», изменение размера изображения) через генерацию HDL-кода в Simulink и вручную написанный Verilog. Оптимизация задержки и пропускной способности, решение проблем синхронизации и конвейеризации.}
            \resumeItem{\textbf{Симуляция, верификация и синтез} \\ Создание всесторонних тестовых стендов на Verilog и использование Python для автоматизации симуляции и анализа данных для проверки функциональности и производительности DSP. Синтез, размещение и трассировка HDL-кода с использованием Vivado (FPGA: Xilinx Artix-7, Zybo Z7). Верификация алгоритмов ISP путем передачи тестовых изображений через HDMI с хост-компьютера, а затем с живой камеры, подключенной к FPGA.}
           \resumeItem{\textbf{Исследование устройств FPGA под воздействием электронно-плазменного излучения} \\ Проведение экспериментов по облучению плат FPGA электронными пучками (25 -- 60 кэВ, до 100 мА) в низкодавленном кислороде ($10^{-6}$ -- 50 Торр), генерируя плазму и рентгеновские лучи. Применение комбинированного термоциклирования (218--393 K) и поверхностного заряда для оценки надежности FPGA под радиационным и плазменным стрессом.}

        \resumeItemListEnd

    %--- Previous (concurrent) roles ---
    \resumeSubheading
        {\href{https://new.ras.ru/en/}{\parbox[t]{\dimexpr\linewidth-4cm}{Институт проблем технологии микроэлектроники и особочистых материалов Российской академии наук}}}{Март 2023 -- Август 2024}
        {\textbf{Лабораторный исследователь}}{Москва, Россия}
        \resumeItemListStart
            \resumeItem{\textbf{Установка и автоматизация экспериментального оборудования} \\ Установка экспериментальных устройств, включая Everbeing Cryo-station (80K – 450K) с 4 микроманипуляторами, Lakeshore Temperature Controller Model 336, Keithley SourceMeter 2450, Parametric Analyzer Keithley 4200A-SCS, Keysight Electrometer B2987A, Aktakom 3048 и Zurich Instruments MFIA Impedance Analyzer. Разработка {\href{https://github.com/rezaaliasgarirenani/IMT-Automation}{приложений}} в MATLAB для автоматизации экспериментальных методов: термостимулированный ток, емкость-напряжение, ток-напряжение, ток-время и спектроскопия глубоких уровней с переходными процессами.}
            \resumeItem{\textbf{Теоретическое исследование} \\ Разработка теоретического понимания и изучение экспериментальных методов для полупроводниковых устройств (MOS, MOSFET, диод, RRAM). Прогнозирование поведения образцов и определение параметров измерений.}
            \resumeItem{\textbf{Экспериментальное исследование и обработка данных} \\ Проведение экспериментов по электрической характеризации микроэлектронных структур, обработка данных с использованием MATLAB и Origin Pro, удаление посторонних точек случайного телеграфного шума, сравнение результатов с теоретическими моделями.}
        \resumeItemListEnd
    \resumeSubheading
        {\href{https://mipt.ru/dasr/about/kaf_faculty/mmsp}{\parbox[t]{\dimexpr\linewidth-6cm}{Лаборатория моделирования механических систем и процессов}}}{Март 2023 -- Август 2024}
        {\textbf{Инженер / Техник}}{Москва, Россия}
        \resumeItemListStart
            \resumeItem{\textbf{Инженерное проектирование и разработка} \\ Проектирование, разработка и анализ моделей для орбитального развертывателя CubeSat и вибрационного крепления с использованием SolidWorks. Создание нескольких прототипов, прошедших симуляцию случайных вибраций и динамический анализ, тестирование на вибростенде UVE 4000 для механических факторов окружающей среды и устойчивости к вибрациям.}
        \resumeItemListEnd
        
\resumeSubHeadingListEnd
\vspace{-10pt} % Adjusted spacing

%-----------PUBLICATIONS---------------
\section{\textcolor{red}{Пу}бликации и конференции}
\resumeSubHeadingListStart
    \item \underline{Р. Алиасгари Ренани}, О.А. Солтанович, М.А. Князев, С.В. Ковешников. \\ \textit{Исследование устройств MOS на основе \ce{SiO2}, облученных низкоэнергетичными электронами, методами емкость-напряжение и термостимулированный ток.} \href{https://doi.org/10.1134/S1063739723600516}{Журнальная статья}, Russian Microelectronics, 2023
    \item \underline{Р. Алиасгари Ренани}, О.А. Солтанович, М.А. Князев, С.В. Ковешников. \\ \textit{Изучение устройств MOS на основе \ce{SiO2} методами емкость-напряжение и термостимулированный ток.} \href{https://icmne.ftian.ru/wp-content/uploads/icmne-2023_e-version.pdf}{Презентация}, p.122. 15-я Международная конференция по микро- и наноэлектронике (\href{https://icmne.ftian.ru}{ICMNE 2023})
    \item \underline{Р. Алиасгари Ренани}, О.А. Солтанович, М.А. Князев, С.В. Ковешников. \\ \textit{Исследование электрически активных дефектов, введенных в оксид кремния облучением низкоэнергетичными электронами, методами характеристик емкость-напряжение и термостимулированного тока.} \href{https://cebt23.iptm.ru/download/numbered/91.pdf}{Постер}, Вторая совместная конференция по технологиям электронных пучков и рентгеновской оптике в микроэлектронике (\href{https://cebt23.iptm.ru}{CALT 2023})
    \item \underline{В. Вологин}, Р. Алиасгари Ренани, В. Васильевский, В. Чесноков \\ \textit{Сравнительный анализ ручного Verilog и HDL-кода, сгенерированного в Simulink, для алгоритмов обработки изображений.} Предстоящая презентация, \href{https://meetups.yadro.com/fpga-msk-1125/}{Конференция FPGA-Systems 2025}
    
\resumeSubHeadingListEnd
\vspace{-12pt} % Adjusted spacing

%-----------TECHNICAL SKILLS-----------
\section{\textcolor{red}{Те}хнические навыки}
\begin{center}
    \begin{minipage}[t]{0.45\textwidth}
        \textbf{Продвинутый уровень}
        \vspace{-7pt}
        \begin{itemize}
            \setlength{\itemindent}{0em}
            \setlength{\itemsep}{-3pt}
            \item Экспериментальная автоматизация
            \item Электрическая характеризация
            \item Обработка данных
            \item MATLAB, Simulink, HDL Coder
            \item Verilog, проектирование RTL
            \item Разработка FPGA
            \item Вычисления с фиксированной точкой
            \item SciPy, NumPy, OpenCV
        \end{itemize}
    \end{minipage}%
    \hfill
    \begin{minipage}[t]{0.45\textwidth}
        \textbf{Средний уровень}
        \vspace{-7pt}
        \begin{itemize}
            \setlength{\itemindent}{0em}
            \setlength{\itemsep}{-3pt}
            \item Python, C++, Arduino
            \item Vivado, Vitis
            \item Git, Unix/Linux OS
            \item OriginLab
            \item SolidWorks
            \item ERDAS IMAGINE
            \item PCB, EasyEDA
            \item OpenRocket
        \end{itemize}
    \end{minipage}
\end{center}
\vspace{-10pt} % Adjusted spacing

%-----------PROJECTS-----------
\section{\textcolor{red}{Пр}оекты}
    %--- Most recent projects first ---
    {\href{https://github.com/icarus-imperium/rocket-2025}{\textbf{Модельная ракета}}} \hfill \textbf{Апрель 2023 и 2025} \\
    \begin{itemize}
        \vspace{-9pt}
    \item Построил модельные ракеты с импульсами 40 и 60 Ньютон-секунд в составе команды.
        \vspace{-9pt}
    \item Запустил три модельные ракеты за два года на День космонавтики МФТИ.
    \end{itemize}
    \vspace{-9pt}
    {\href{https://github.com/rezaaliasgarirenani/Rover}{\textbf{Модель лунохода}}} \hfill \textbf{Июнь 2023}
    \begin{itemize}
        \vspace{-9pt}
    \item Сотрудничал над машиной, способной преодолевать препятствия без круглых колес.
        \vspace{-9pt}
    \item Протестировал несколько прототипов, финальный дизайн одобрен главой лаборатории.
    \end{itemize}
    \vspace{-9pt}
    {\href{https://github.com/rezaaliasgarirenani/Aircraft-Detection-System}{\textbf{Система обнаружения самолетов}}} \hfill \textbf{Июнь 2022}
    \begin{itemize}
        \vspace{-9pt}
    \item Исследовал и применил алгоритмы для обнаружения самолетов с использованием фоторезисторов и транзисторов.
        \vspace{-9pt}
    \item Разработал систему, способную к ротационному и трансляционному движению для отслеживания самолетов.
    \end{itemize}
    \vspace{-9pt}
    {\href{https://github.com/rezaaliasgarirenani/Non-Conservative-Electric-Fields-and-Voltmeters}{\textbf{Исследование неконсервативных электрических полей и вольтметров}}} \hfill \textbf{Май 2022} \\
    \begin{itemize}
        \vspace{-9pt}
    \item Разработал экспериментальную установку для анализа зависимости показаний вольтметра от позиции в параллельных цепях.
        \vspace{-9pt}
    \item Продемонстрировал неинтуитивные разности потенциалов, генерируемые изменяющимися магнитными полями.
    \end{itemize}
    \vspace{-10pt} % Adjusted spacing

%-----------AWARDS-----------
\section{\textcolor{red}{На}грады}
\begin{itemize}
    %--- Most recent award first ---
     \resumeItem{Полная государственная российская стипендия для иностранных студентов, МФТИ} \hfill \textbf{Сентябрь 2024}
    \resumeItem{Иранская государственная стипендия, Исфаханский технологический университет} \hfill \textbf{Сентябрь 2019}
    \resumeItem{Участник \href{https://congress.aero/en/}{5-го и 7-го Евразийского аэрокосмического конгресса}} \hfill \textbf{Июль 2023 и 2025}

\end{itemize}
\vspace{-10pt} % Adjusted spacing

%-----------LANGUAGES-----------
\section{\textcolor{red}{Яз}ыки}
\begin{center}
    \textbf{Английский}: C2 (TOEFL iBT 113) \quad $|$ \quad 
    \textbf{Русский}: B1-B2 \quad $|$ \quad 
    \textbf{Персидский}: Родной
\end{center}

\end{document}